\subsection{Haces de tipo finito}
\label{subsection:0.5.2-fr}

\begin{env}[5.2.1]
\label{0.5.2.1-fr}
Se dice que un $\mathcal{O}_X$-Módulo $\mathcal{F}$ es \emph{de tipo finito} si,
para todo $x\in X$, existe una vecindad abierta $U$ de $x$ tal que
$\mathcal{F}|U$ está generado por una familia \emph{finita} de secciones
sobre $U$, o, equivalentemente, es isomorfo a un haz cociente de un
haz de la forma $(\mathcal{O}_X|U)^p$, donde $p$ es finito. Todo haz
cociente de un haz de tipo finito es de tipo finito, así como toda suma
directa finita y todo producto tensorial finito de haces de tipo finito. Un
$\mathcal{O}_X$-Módulo de tipo finito no es necesariamente cuasicoherente,
como se puede ver para el $\mathcal{O}_X$-Módulo
$\mathcal{O}_X/\mathcal{F}$, donde $\mathcal{F}$ es el ejemplo de (5.1.1). Si
$\mathcal{F}$ es de tipo finito, $\mathcal{F}_x$ es un
$\mathcal{O}_x$-módulo de tipo finito para todo $x\in X$, pero el ejemplo
(5.1.1) muestra que esta condición necesaria no es en general
suficiente.
\end{env}

\begin{env}[5.2.2]
\label{0.5.2.2-fr}
Sea $\mathcal{F}$ un $\mathcal{O}_X$-Módulo \emph{de tipo finito}. Si $s_i$
$(1\leq i\leq n)$ son secciones de $\mathcal{F}$ sobre una vecindad
abierta $U$ de un punto $x\in X$ y si los $(s_i)_x$ generan
$\mathcal{F}_x$, existe una vecindad abierta $V\subset U$ de $x$ tal que
los $(s_i)_y$ generan $\mathcal{F}_y$ para todo $y\in V$ (FAC, I, 2, 12,
prop.~1). Se concluye, en particular, que el soporte de $\mathcal{F}$ es
\emph{cerrado}.

\oldpage[0\textsubscript{I}]{46}

Del mismo modo, si $u:\mathcal{F}\to\mathcal{G}$ es un homomorfismo tal que
$u_x=0$, entonces existe una vecindad $U$ de $x$ tal que $u_y=0$ para todo
$y\in U$.
\end{env}

\begin{env}[5.2.3]
\label{0.5.2.3-fr}
Supongamos que $X$ es \emph{cuasicompacto}, y sean $\mathcal{F}$, $\mathcal{G}$ dos
$\mathcal{O}_X$-Módulos tales que $\mathcal{G}$ es de \emph{tipo finito},
$u:\mathcal{F}\to\mathcal{G}$ un homomorfismo \emph{sobreyectivo}. Supongamos
además que $\mathcal{F}$ es límite inductivo de un sistema inductivo
$(\mathcal{F}_\lambda)$ de $\mathcal{O}_X$-Módulos. Entonces existe un índice
$\mu$ tal que el homomorfismo $\mathcal{F}_\mu\to\mathcal{G}$ es
\emph{sobreyectivo}. En efecto, para todo $x\in X$, existe un sistema finito de
secciones $s_i$ de $\mathcal{G}$ sobre una vecindad abierta $U(x)$ de $x$
tal que los $(s_i)_y$ generan $\mathcal{G}_y$ para todo $y\in U(x)$; hay,
pues, una vecindad abierta $V(x)\subset U(x)$ de $x$ y $n$ secciones $t_i$ de
$\mathcal{F}$ sobre $V(x)$ tales que
$s_i|V(x)=u(t_i)$ para todo $i$; se puede suponer además que los $t_i$ son
las imágenes canónicas de secciones de un mismo haz
$\mathcal{F}_{\lambda(x)}$ sobre $V(x)$. Cubramos entonces $X$ con un
número finito de vecindades $V(x_k)$, y sea $\mu$ un índice superior a todos
los $\lambda(x_k)$; es claro que este índice responde a la cuestión.

Suponiendo todavía que $X$ es cuasicompacto, sea $\mathcal{F}$ un
$\mathcal{O}_X$-Módulo de tipo finito generado por sus secciones sobre $X$
(5.1.1); entonces $\mathcal{F}$ está generado por una subfamilia \emph{finita}
de estas secciones: basta, en efecto, cubrir $X$ con un número finito de
vecindades abiertas $U_k$ tales que, para cada $k$, hay un número finito de
secciones $s_{ik}$ de $\mathcal{F}$ sobre $X$ cuyas restricciones a
$U_k$ generan $\mathcal{F}|U_k$; es claro que los $s_{ik}$ generan
entonces $\mathcal{F}$.
\end{env}

\begin{env}[5.2.4]
\label{0.5.2.4-fr}
Sea $f:X\to Y$ un morfismo de espacios anillados. Si $\mathcal{G}$ es un
$\mathcal{O}_Y$-Módulo de tipo finito, entonces $f^*(\mathcal{G})$ es un
$\mathcal{O}_X$-módulo de tipo finito. En efecto, para todo $x\in X$, hay una
vecindad abierta $V$ de $f(x)$ en $Y$ y un homomorfismo sobreyectivo
$v:\mathcal{O}_Y^p|V\to\mathcal{G}|V$. Si $U=f^{-1}(V)$ y si $f_U$ es la
restricción de $f$ a $U$, se tiene
$f^*(\mathcal{G})|U=f_U^*(\mathcal{G}|V)$; como $f_U^*$ es exacto por la derecha
(4.3.1) y conmuta con las sumas directas (4.3.2), $f_U^*(v)$ es un
homomorfismo sobreyectivo
$\mathcal{O}_X^p|U\to f^*(\mathcal{G})|U$.
\end{env}

\begin{env}[5.2.5]
\label{0.5.2.5-fr}
Se dice que un $\mathcal{O}_X$-Módulo $\mathcal{F}$ \emph{admite una presentación
finita} si, para todo $x\in X$, existe una vecindad abierta $U$ de $x$ tal
que $\mathcal{F}|U$ es isomorfo a un \emph{conúcleo de un
$(\mathcal{O}_X|U)$-homomorfismo}
$\mathcal{O}_X^p|U\to\mathcal{O}_X^q|U$, siendo $p$ y $q$ dos enteros $>0$.
Tal $\mathcal{O}_X$-Módulo es, por tanto, de tipo finito y cuasicoherente. Si
$f:X\to Y$ es un morfismo de espacios anillados, y si $\mathcal{G}$ es un
$\mathcal{O}_Y$-Módulo que admite una presentación finita,
$f^*(\mathcal{G})$ admite una presentación finita, como muestra el
razonamiento de (5.1.4).
\end{env}

\begin{env}[5.2.6]
\label{0.5.2.6-fr}
Sea $\mathcal{F}$ un $\mathcal{O}_X$-Módulo que admite una presentación finita
(5.2.5); entonces, para todo $\mathcal{O}_X$-Módulo $\mathcal{H}$,
el homomorfismo canónico funtorial
\[
  \bigl(\mathcal{H}\!om_{\mathcal{O}_X}(\mathcal{F},\mathcal{H})\bigr)_x
  \longrightarrow
  \operatorname{Hom}_{\mathcal{O}_x}(\mathcal{F}_x,\mathcal{H}_x)
\]
es \emph{biyectivo} (T, 4.1.1).
\end{env}

\begin{env}[5.2.7]
\label{0.5.2.7-fr}
Sean $\mathcal{F}$, $\mathcal{G}$ dos $\mathcal{O}_X$-Módulos que tienen una
presentación finita. Si, para un $x\in X$, $\mathcal{F}_x$ y $\mathcal{G}_x$
son dos $\mathcal{O}_x$-módulos \emph{isomorfos}, entonces existe una
vecindad abierta $U$ de $x$ tal que $\mathcal{F}|U$ y $\mathcal{G}|U$ son
\emph{isomorfos}. En efecto, si
$\varphi:\mathcal{F}_x\to\mathcal{G}_x$ y
$\psi:\mathcal{G}_x\to\mathcal{F}_x$ son dos isomorfismos inversos, existe,
según (5.2.6), una vecindad abierta $V$ de $x$ y una sección $u$
(resp. $v$) de
$\mathcal{H}\!om_{\mathcal{O}_X}(\mathcal{F},\mathcal{G})$ (resp.
$\mathcal{H}\!om_{\mathcal{O}_X}(\mathcal{G},\mathcal{F})$) sobre $V$
tal

\oldpage[0\textsubscript{I}]{47}

que $u_x=\varphi$ (resp. $v_x=\psi$). Como $(u\circ v)_x$ y
$(v\circ u)_x$ son los automorfismos identidad, existe una vecindad
abierta $U\subset V$ de $x$ tal que $(u\circ v)|U$ y $(v\circ u)|U$ son los
automorfismos identidad, de donde la proposición.
\end{env}

